% --- THE MAGIC TOGGLE ---
% STUDENT VERSION: \documentclass[12pt, addpoints]{exam}
% TEACHER VERSION: \documentclass[12pt, answers, addpoints]{exam}
\documentclass[12pt, answers, addpoints]{exam}

\usepackage{amsmath}
\usepackage{listings} % For formatting code snippets
\usepackage{xcolor}   % For styling code
\usepackage{geometry}

\geometry{margin=1in}

% --- Code Snippet Styling ---
\lstset{
    basicstyle=\ttfamily\small,
    keywordstyle=\color{blue}\bfseries,
    commentstyle=\color{gray},
    stringstyle=\color{red},
    frame=single,
    showstringspaces=false
}

% --- Header and Footer Setup ---
\pagestyle{headandfoot}
\firstpageheader{CS 301: Advanced Algorithms}{}{Midterm Examination}
\runningheader{CS 301}{}{\studentname}
\runningheadrule
\firstpagefooter{}{Page \thepage\ of \numpages}{}
\runningfooter{}{Page \thepage\ of \numpages}{}

\begin{document}

% --- Title Block ---
\begin{center}
    {\huge \textbf{Computer Science Midterm}} \\[1em]
    {\large \textbf{Fall Semester 2026}} \\[2em]
\end{center}

\makebox[\textwidth]{Name:\enspace\hrulefill} \\[1em]
\makebox[\textwidth]{Student ID:\enspace\hrulefill}

\vspace{2em}

% --- Instructions and Grading Table ---
\begin{center}
    \fbox{\fbox{\parbox{5.5in}{\centering
        \textbf{Instructions:} Answer all questions in the space provided. Show your work for partial credit. 
        Calculators are not permitted. You have 90 minutes to complete this exam.
    }}}
\end{center}

\vspace{1em}

% This table automatically calculates total points based on the \question[points] tags!
\begin{center}
    \gradetable[h][questions]
\end{center}

\vspace{2em}
\hrule
\vspace{2em}

% --- The Questions ---
\begin{questions}

\section*{Part 1: Multiple Choice}

\question[5] What is the worst-case time complexity of a standard quicksort algorithm?
\begin{choices}
    \choice $O(n)$
    \choice $O(n \log n)$
    \CorrectChoice $O(n^2)$
    \choice $O(1)$
\end{choices}

\question[5] Which of the following HTTP methods is intended to be idempotent?
\begin{checkboxes}
    \choice POST
    \CorrectChoice PUT
    \choice PATCH
    \choice None of the above
\end{checkboxes}


\section*{Part 2: Short Answer}

\question[10] Explain the difference between a stack and a queue. Provide one practical use case for each in software engineering.
% \fillwithlines creates blank lines for the student, but gets replaced by the solution in the answer key.
\begin{solution}[2.5in]
    A stack operates on a Last-In, First-Out (LIFO) principle. A practical use case is managing the "Undo" feature in a text editor or tracking the call stack in memory.\\
    
    A queue operates on a First-In, First-Out (FIFO) principle. A practical use case is managing background tasks in a message broker or handling requests to a shared printer.
\end{solution}

\section*{Part 3: Code Analysis}

\question[15] Identify the bug in the following Python function designed to check if a string is a palindrome.

\begin{lstlisting}[language=Python]
def is_palindrome(text):
    text = text.lower()
    reversed_text = text[::-1]
    if text = reversed_text:
        return True
    else:
        return False
\end{lstlisting}

\begin{solution}[1.5in]
    The bug is on line 4: \texttt{if text = reversed\_text:}. \\
    This is an assignment operator, but it should be an equality operator (\texttt{==}). The corrected line should read: \texttt{if text == reversed\_text:}.
\end{solution}


\section*{Part 4: Mathematical Analysis}

\question[15] Use mathematical induction to prove that the sum of the first $n$ positive integers is given by the formula:
$$ \sum_{i=1}^{n} i = \frac{n(n+1)}{2} $$

% \makeemptybox creates a blank square for students to do math in.
\begin{solution}[3.5in]
    \textbf{Base Case ($n=1$):} \\
    LHS: $\sum_{i=1}^{1} i = 1$ \\
    RHS: $\frac{1(1+1)}{2} = \frac{2}{2} = 1$. The base case holds. \\[1em]
    \textbf{Inductive Step:} \\
    Assume true for $n=k$: $\sum_{i=1}^{k} i = \frac{k(k+1)}{2}$. \\
    We must show it holds for $n=k+1$: $\sum_{i=1}^{k+1} i = \frac{(k+1)(k+2)}{2}$. \\
    $\sum_{i=1}^{k+1} i = \left(\sum_{i=1}^{k} i\right) + (k+1)$ \\
    Substitute the inductive hypothesis: \\
    $= \frac{k(k+1)}{2} + (k+1)$ \\
    $= \frac{k(k+1) + 2(k+1)}{2} = \frac{(k+1)(k+2)}{2}$. \\
    Thus, by mathematical induction, the statement is true for all $n \geq 1$.
\end{solution}

\end{questions}

\end{document}